% Part: model-theory
% Chapter: basics
% Section: theory-of-m

\documentclass[../../../include/open-logic-section]{subfiles}

\begin{document}

\olfileid{mod}{bas}{thm}

\section{نظریهٔ یک \printtoken{S}{structure}}

هر !!{structure}~$\Struct{M}$ برخی !!{sentence}s را صادق و برخی را
کاذب می‌سازد. مجموعهٔ همهٔ !!{sentence}sی که آن ساختار صادق می‌سازد
\emph{نظریهٔ} آن نامیده می‌شود. این مجموعه در واقع یک نظریه است، زیرا
هرچه از آن نتیجه شود باید در همهٔ مدل‌هایش، از جمله~$\Struct{M}$، صادق
باشد.

\begin{defn}
  به ازای !!a{structure}~$\Struct M$، \emph{نظریهٔ} $\Struct{M}$
  مجموعهٔ~$\Theory{M}$ از !!{sentence}sی است که در~$\Struct{M}$ صادق‌اند؛
  یعنی $\Theory{M} = \Setabs{!A}{\Sat{M}{!A}}$.
\end{defn}

اصطلاح «نظریه» را همچنین به‌طور غیررسمی برای اشاره به مجموعه‌هایی از
!!{sentence}s دارای تعبیر مقصود به کار می‌بریم، خواه از نظر استنتاجی
بسته باشند خواه نباشند.

\begin{prop}
برای هر $\Struct{M}$، $\Theory{M}$ کامل است.
\end{prop}

\begin{proof}
برای هر !!{sentence}~$!A$، یا $\Sat{M}{!A}$ یا
$\Sat{M}{\lnot !A}$؛ پس یا $!A \in \Theory{M}$ یا
$\lnot !A \in \Theory{M}$.
\end{proof}

\begin{prop}\ollabel{prop:equiv}
  اگر برای هر $!A \in \Theory{M}$، $\Struct{N} \models !A$، آنگاه
  $\Struct{M} \elemequiv \Struct{N}$.
\end{prop}

\begin{proof}
چون برای همهٔ $!A \in \Theory{M}$، $\Sat{N}{!A}$، داریم
$\Theory{M} \subseteq \Theory{N}$. اگر $\Sat{N}{!A}$، آنگاه
$\Sat/{N}{\lnot !A}$، پس $\lnot !A \notin \Theory{M}$. چون
$\Theory{M}$ کامل است، $!A \in \Theory{M}$. بنابراین
$\Theory{N} \subseteq \Theory{M}$، و داریم
$\Struct{M} \elemequiv \Struct{N}$.
\end{proof}

\begin{rem}\ollabel{remark:R}
  $\Struct{R} = \langle\Real, <\rangle$ را در نظر بگیرید؛ یعنی
  !!{structure}ای که دامنه‌اش مجموعهٔ~$\Real$ از اعداد حقیقی است، در
  !!{language}ی که تنها از یک !!{predicate} دوموضعی تشکیل می‌شود که به
  رابطهٔ~$<$ روی اعداد حقیقی تعبیر شده است. آشکارا $\Struct{R}$
  !!{nonenumerable} است؛ اما چون $\Theory{R}$ آشکارا سازگار است، بنا بر
  قضیهٔ لوونهایم--اسکولم !!a{enumerable} مدل، مثلاً $\Struct{S}$، دارد
  و بنا بر \olref{prop:equiv}، $\Struct{R} \equiv \Struct{S}$. افزون
  بر این، چون $\Struct{R}$ و~$\Struct{S}$ همریخت نیستند، این نشان
  می‌دهد که عکس \olref[iso]{thm:isom} به‌طور کلی برقرار نیست.
\end{rem}

\end{document}
